%------------------------------------------------------------------------------%
\addtocounter{exemplo}{1}
\begin{frame}{Exemplo \theexemplo}
  Encontre o ponto do plano \(x+2y+3z=13\) mais próximo do ponto \((1, 1, 1)\)
\end{frame}

%------------------------------------------------------------------------------%
\begin{frame}{Exemplo \theexemplo}
  Queremos minimizar a distância
  \[
    d(x, y, z) = \sqrt{(x-1)^2 + (y-1)^2 + (z-1)^2}
  \]
  Como a função raiz quadrada é crescente podemos minimizar a função
  \[
    f(x, y, z) = (x-1)^2 + (y-1)^2 + (z-1)^2
  \]
  suas derivadas parciais são
  \[
    f_x(x, y, z) = 2(x-1)
    \qquad
    f_y(x, y, z) = 2(y-1)
    \qquad
    f_z(x, y, z) = 2(z-1)
  \]
\end{frame}

%------------------------------------------------------------------------------%
\begin{frame}{Exemplo \theexemplo}
  As derivadas parciais da função \(g(x, y, z)=x+2y+3z\) são
  \[
    g_x(x, y, z) = 1
    \qquad
    g_y(x, y, z) = 2
    \qquad
    g_z(x, y, z) = 3
  \]
\end{frame}

%------------------------------------------------------------------------------%
\begin{frame}{Exemplo \theexemplo}
  Aplicando os Multiplicadores de Lagrange, precisamos resolver o sistema
  \[
  \begin{cases}
    2(x-1) =  \lambda \\
    2(y-1) = 2\lambda \\
    2(z-1) = 3\lambda \\
    x + 2y + 3z = 13
  \end{cases}
  \]
\end{frame}

%------------------------------------------------------------------------------%
\begin{frame}{Exemplo \theexemplo}
  Das três primeiras equações obtemos
  \[
    2(x-1) = \lambda
    \quad\Rightarrow\quad
    x = \frac{\lambda}{2} + 1
  \]
  \[
    2(y-1) = 2\lambda
    \quad\Rightarrow\quad
    y = \lambda + 1
  \]
  \[
    2(z-1) = 3\lambda
    \quad\Rightarrow\quad
    z = \frac{3\lambda}{2} + 1
  \]
\end{frame}

%------------------------------------------------------------------------------%
\begin{frame}{Exemplo \theexemplo}
  Substituindo na última equação
  \begin{align*}
    x + 2y + 3z & = 13
    \\[2mm]
    \left(\frac{\lambda}{2} + 1\right)
    + 2 \left(\lambda + 1\right)
    + 3 \left(\frac{3\lambda}{2} + 1\right) & = 13
    \\[2mm]
    \frac{\lambda}{2} + 1
    + 2\lambda + 2
    + \frac{9\lambda}{2} + 3 & = 13
    \\[2mm]
    \lambda + 2
    + 4\lambda + 4
    + 9\lambda + 6 & = 13\times 2
    \\[2mm]
    14\lambda + 12 & = 26
    \\[2mm]
    14\lambda & = 14
    \\[2mm]
    \lambda & = 1
  \end{align*}
\end{frame}

%------------------------------------------------------------------------------%
\begin{frame}{Exemplo \theexemplo}
  Portanto
  \[
    x = \frac{\lambda}{2} + 1 = \frac{3}{2}
  \]
  \[
    y = \lambda + 1 = 2
  \]
  \[
    z = \frac{3\lambda}{2} + 1 = \frac{5}{2}
  \]
  O ponto mais próximo é o ponto \(\left(\frac{3}{2}, 2, \frac{5}{2}\right)\)
\end{frame}

%------------------------------------------------------------------------------%
